% \hold{...} underlines its argument (a syllable held on a longer note).  The
% rule sits a fixed \holddepth below the baseline -- just under the descenders,
% so underlines line up across syllables -- and is \holdthick thick (set
% independent of the mark size).  It is also \holdgap narrower than the syllable
% and centered, so the underlines of two adjacent held syllables stay separate
% rather than joining.  These dimensions are set from the font in layout.tex.
% The \rlap overlays the rule without consuming width; \leavevmode lets \hold
% start a line.
\newdimen\holddepth
\newdimen\holdthick
\newdimen\holdgap
\def\hold#1{%
  \leavevmode
  \setbox0=\hbox{#1}%
  \dimen0=\wd0 \advance\dimen0 by-\holdgap   % shorten the rule by a fixed amount
  \rlap{\kern.5\holdgap                       % ... and centre it under the syllable,
    \lower\holddepth\hbox{\vrule height0pt depth\holdthick width\dimen0}}%
  \box0                                        % leaving half the gap on each side
}
